Die politische Linke in Deutschland steht vor einer historischen Aufgabe, die über parlamentarische Routine hinausgeht. Die vorliegende Analyse hat gezeigt, dass die zentralen Begriffe unserer Zeit – Faschismus, Nationalsozialismus, Judenhass, Zionismus und Sozialismus – systematisch verdreht werden. Diese Verdrehung ist keine akademische Nebensächlichkeit, sondern ein zentrales Instrument der Herrschaftssicherung.

Die Konsequenzen für die politische Arbeit der Linken sind weitreichend:

\begin{enumerate}
    \item \textbf{Begriffliche Klarheit schaffen:} Die Linke muss die Begriffe zurückerobern, die ihr geraubt wurden. Faschismus darf nicht länger mit Nationalsozialismus gleichgesetzt werden. Der völkische Kampfbegriff \emph{Antisemitismus} darf nicht länger gegen Zionismuskritik missbraucht werden. Sozialismus darf nicht länger auf Stalin, Honecker oder Mao reduziert werden.
    
    \item \textbf{Die CDU/CSU als faschistische Partei benennen:} Die strukturelle Definition des Faschismus – Steuerung der Politik durch das Kapital, antidemokratisch in der Wirkung – trifft auf die CDU/CSU zu. Diese Benennung ist keine Polemik, sondern das Ergebnis einer historisch einwandfreien Analyse.
    
    \item \textbf{Die Instrumentalisierung des Judenhass-Begriffs entlarven:} Der reale Judenhass muss bekämpft werden. Aber der Missbrauch des völkisch geprägten Begriffs \emph{Antisemitismus} zur Delegitimierung von Antizionismus muss benannt und zurückgewiesen werden. Die Linke verurteilt Judenhass – aber sie verurteilt auch den Zionismus.
    
    \item \textbf{Internationale Solidarität praktisch werden lassen:} Die Linke muss sich solidarisch zeigen mit allen Völkern und Bewegungen, die unter imperialer Politik, kolonialer Unterdrückung und kapitalistischer Ausbeutung leiden – insbesondere mit den Palästinensern.
    
    \item \textbf{Den Sozialismus als demokratische Alternative verteidigen:} Die historischen Verirrungen des Sozialismus (Stalin, Honecker, Mao) müssen kritisch aufgearbeitet werden. Aber sie dürfen nicht dazu führen, dass das sozialistische Ideal selbst aufgegeben wird – die Befreiung der Arbeiterklasse vom Kapital, internationale Solidarität, echte Demokratie.
\end{enumerate}

Die Analyse zeigt: Der Faschismus sitzt nicht nur am rechten Rand. Er sitzt in der Mitte – und er trägt Anzug, nicht Uniform. Aber er ist nicht weniger gefährlich. Die CDU/CSU, die Trump und Netanjahu unterstützt, die Israel bedingungslos verteidigt und jede Kritik daran mit dem völkischen Kampfbegriff \emph{Antisemitismus} diffamiert – indem sie Antizionismus fälschlicherweise mit Judenhass gleichsetzt – ist die gegenwärtige Erscheinungsform des Faschismus in Deutschland.

Es ist an der Zeit, die Begriffe nicht länger verdrehen zu lassen. Es ist an der Zeit, den Kampf gegen den Faschismus in der Mitte aufzunehmen – und nicht nur am Rand.